\subsubsection*{Argument Forms}

\begin{remark*}[Working vocabulary]
This topic introduces argument forms, premises, conclusions, substitution
instances, valid argument forms, invalid argument forms, and counterexample
truth assignments.
\end{remark*}

% BEGIN GENERATED ARTIFACT: def:argument-form-propositional-logic
\begin{definitionbox}{Definition (Argument Form)}
\begin{definition}[Argument Form]
\label{def:argument-form-propositional-logic}
An \emph{argument form} is a schematic pattern consisting of a finite list of
premise formulas and one conclusion formula.
\end{definition}
\end{definitionbox}

\begin{remark*}[Standard quantified statement]
\[
\operatorname{ArgumentForm}(\Gamma,\varphi)
\Longleftrightarrow
\Gamma\subseteq\WFF\text{ is finite and }\varphi\in\WFF.
\]
\end{remark*}

\begin{remark*}[Definition predicate reading]
\(\operatorname{ArgumentForm}(\Gamma,\varphi)\) means that \(\Gamma\) is the set
or finite list of premises and \(\varphi\) is the conclusion.
\end{remark*}

\begin{remark*}[Interpretation]
An argument form records a pattern of inference independently of the particular
subject matter. Its validity depends on truth preservation, not on whether the
individual formulas happen to be true in one example.
\end{remark*}

\begin{dependencies}
\begin{itemize}
  \item \hyperref[def:well-formed-formula]{Well-Formed Formula}
\end{itemize}
\end{dependencies}
% END GENERATED ARTIFACT: def:argument-form-propositional-logic

% BEGIN GENERATED ARTIFACT: def:valid-argument-form-propositional-logic
\begin{definitionbox}{Definition (Valid Argument Form)}
\begin{definition}[Valid Argument Form]
\label{def:valid-argument-form-propositional-logic}
An argument form with premises \(\Gamma\) and conclusion \(\varphi\) is
\emph{valid} if every truth assignment that satisfies every formula in
\(\Gamma\) also satisfies \(\varphi\).
\end{definition}
\end{definitionbox}

\begin{remark*}[Standard quantified statement]
\[
\operatorname{ValidArgumentForm}(\Gamma,\varphi)
\Longleftrightarrow
(\forall v)\bigl((\forall\gamma\in\Gamma)(v\models\gamma)
\to v\models\varphi\bigr).
\]
\end{remark*}

\begin{remark*}[Definition predicate reading]
\(\operatorname{ValidArgumentForm}(\Gamma,\varphi)\) means that truth of all
premises forces truth of the conclusion under every truth assignment.
\end{remark*}

\begin{remark*}[Interpretation]
Validity is a preservation property. A valid argument form may have false
premises in a particular instance, but it cannot have true premises and a false
conclusion under the same truth assignment.
\end{remark*}

\begin{dependencies}
\begin{itemize}
  \item \hyperref[def:argument-form-propositional-logic]{Argument Form}
  \item \hyperref[def:satisfaction-propositional-formula]{Satisfaction of a Formula}
  \item \hyperref[def:logical-consequence-propositional-logic]{Logical Consequence}
\end{itemize}
\end{dependencies}
% END GENERATED ARTIFACT: def:valid-argument-form-propositional-logic

% BEGIN GENERATED ARTIFACT: def:counterexample-assignment-propositional-logic
\begin{definitionbox}{Definition (Counterexample Assignment)}
\begin{definition}[Counterexample Assignment]
\label{def:counterexample-assignment-propositional-logic}
A \emph{counterexample assignment} to an argument form \((\Gamma,\varphi)\) is a
truth assignment that satisfies every premise in \(\Gamma\) and does not satisfy
\(\varphi\).
\end{definition}
\end{definitionbox}

\begin{remark*}[Standard quantified statement]
\[
\operatorname{CounterexampleAssignment}(v,\Gamma,\varphi)
\Longleftrightarrow
(\forall\gamma\in\Gamma)(v\models\gamma)\land v\not\models\varphi.
\]
\end{remark*}

\begin{remark*}[Definition predicate reading]
\(\operatorname{CounterexampleAssignment}(v,\Gamma,\varphi)\) means that \(v\)
shows the premises can all be true while the conclusion is false.
\end{remark*}

\begin{remark*}[Interpretation]
A counterexample assignment is the semantic witness to invalidity. Finding one
settles the question: the argument form is not truth-preserving.
\end{remark*}

\begin{dependencies}
\begin{itemize}
  \item \hyperref[def:valid-argument-form-propositional-logic]{Valid Argument Form}
  \item \hyperref[def:truth-assignment-propositional-logic]{Truth Assignment}
\end{itemize}
\end{dependencies}
% END GENERATED ARTIFACT: def:counterexample-assignment-propositional-logic

% BEGIN GENERATED ARTIFACT: thm:truth-table-validity-test-argument-forms-propositional-logic
\begin{theorembox}{Theorem (Truth-Table Validity Test for Argument Forms)}
\begin{theorem}[Truth-Table Validity Test for Argument Forms]
\label{thm:truth-table-validity-test-argument-forms-propositional-logic}
An argument form \((\Gamma,\varphi)\) is valid if and only if no row of its truth
table makes every formula in \(\Gamma\) true and \(\varphi\) false.

\noindent\hyperref[prf:truth-table-validity-test-argument-forms-propositional-logic]{Go to proof.}
\end{theorem}
\end{theorembox}

\begin{remark*}[Standard quantified statement]
\[
\operatorname{ValidArgumentForm}(\Gamma,\varphi)
\Longleftrightarrow
\neg(\exists v)\operatorname{CounterexampleAssignment}(v,\Gamma,\varphi).
\]
\end{remark*}

\begin{remark*}[Predicate reading]
The argument form is valid exactly when there is no counterexample truth
assignment.
\end{remark*}

\begin{remark*}[Interpretation]
The truth-table test is finite because only the variables appearing in the
premises and conclusion matter. A failed row is a counterexample assignment; the
absence of such rows proves semantic validity.
\end{remark*}

\begin{dependencies}
\begin{itemize}
  \item \hyperref[def:valid-argument-form-propositional-logic]{Valid Argument Form}
  \item \hyperref[def:counterexample-assignment-propositional-logic]{Counterexample Assignment}
  \item \hyperref[def:truth-table-propositional-logic]{Truth Table}
  \item \hyperref[thm:finite-truth-table-propositional-logic]{Finite Truth-Table Theorem}
\end{itemize}
\end{dependencies}
% END GENERATED ARTIFACT: thm:truth-table-validity-test-argument-forms-propositional-logic
